%=======================================================================
% AI Sentinel: Evidence-Grounded Real-Time SOC Platform
% IEEE conference paper base template.
%
% COMPILATION:
%   pdflatex paper.tex
%   pdflatex paper.tex      (run twice for correct references/TOC)
%   bibtex paper            (if using a .bib file; not required here)
%
% Requires: IEEEtran.cls (bundled with every TeX Live / MiKTeX install)
% Figures are referenced but the .png files are optional; the document
% compiles standalone without them.
%=======================================================================
\documentclass[conference]{IEEEtran}

\usepackage{times}                           % Times font
\usepackage{amsmath}
\usepackage{amssymb}
\usepackage{algorithm}
\usepackage{algpseudocode}
\usepackage{graphicx}
\usepackage{booktabs}
\usepackage{array}
\usepackage{multirow}
\usepackage{url}
\usepackage[hidelinks]{hyperref}

\begin{document}

\title{AI Sentinel: An Evidence-Grounded Real-Time Security\\ Operations Center Platform with Provable Data Lineage}

\author{
\IEEEauthorblockN{Akhil A. \thanks{Corresponding author.}}
\IEEEauthorblockA{
Department of Computer Science and Engineering\\
Academic research environment\\
Email: akhil@akhilaids.com\\
}
}

% --------------------------------------------------------------------- %
\begin{abstract}
Security Operations Centers (SOCs) evaluate their effectiveness with
aggregate statistics -- detection coverage, response-time distributions,
and posture scores -- yet most commercially deployed tooling computes these
figures from data whose provenance is unverifiable. Dashboard panels that
report ``100\% coverage'' or ``CRITICAL posture'' are frequently
unreproducible: the underlying counters may be hard-coded, computed over a
window that excludes the relevant observations, or silently substituted with
neutral defaults when a query returns no rows. This paper presents AI
Sentinel, a real-time SOC platform whose central design commitment is
\emph{verifiable ground truth}: every metric the platform reports is derived
by an auditable computation over persisted, tamper-evident records, and
every quantity that cannot be established from real data is returned as an
explicit null or a named absence state rather than a reassuring default.

The system ingests host, network, authentication, and file telemetry through
a non-blocking asynchronous pipeline, applies seventeen runtime-configurable
sliding-window detection rules and a drift-monitored Isolation Forest anomaly
model, correlates detections into MITRE ATT\&CK-mapped incidents, and streams
results to an operator console over an authenticated WebSocket. We contribute
four mechanisms: (i) an HMAC-SHA256 hash-chained audit ledger that makes
administrative history tamper-evident, with verified-chain reporting;
(ii) a *derived* security-posture score in which every point is an explained,
bounded, drill-down delta, and where zero telemetry is scored as the worst
possible state rather than as success; (iii) a *proof-carrying coverage*
analysis that partitions MITRE techniques into
VERIFIED\_LIVE / DETECTING / RULES\_DISABLED / OBSERVED\_NO\_RULES by joining
rule definitions against observed telemetry; and (iv) a data-quality center
that surfaces silent pipeline drops as first-class operational signals.

Evaluation on an 8{,}658-event live corpus shows an ingest-to-console p99 of
10.58~ms (100\% of deliveries under the 2\,s operational target), a
tamper-evident audit ledger of 2{,}847 records verifying intact, and an
alert-compression ratio of 101.7 detection firings per analyst-facing alert.
The evaluation process itself surfaced two silent pipeline defects -- a
NameError that silently dropped an entire malware detection class and a
column allow-list that prevented incident timelines from accumulating -- which
we describe as evidence for the paper's central claim: without a
data-quality signal wired to a live counter, correctness failures of exactly
this kind remain invisible to operators.
\end{abstract}

\begin{IEEEkeywords}
Security operations center, real-time security telemetry, MITRE ATT\&CK,
security posture assessment, tamper-evident logging, evidence integrity,
unsupervised anomaly detection, data provenance, honest computing.
\end{IEEEkeywords}

% --------------------------------------------------------------------- %
\section{Introduction}

\subsection{Problem Statement}

A Security Operations Center exists to answer three questions continuously:
\emph{What is happening now?}, \emph{can we prove what we did about it?}, and
\emph{how good are we at it?}. Commercial SOC platforms answer the third
question with dashboards. The methodological difficulty is that the
measurements behind those dashboards are rarely themselves subjects of
audit. A coverage percentage is usually computed as the ratio of enabled
rules to configured rules -- a statement about the configuration, not about
whether any adversary would actually be detected. A posture score is usually
a weighted sum of dashboard sub-scores, where a component that fails to load
quietly contributes its nominal default. A mean-time-to-detect figure is
frequently averaged over a sample that excludes the very events that were
missed.

The result is a measurement problem that compounds: analysts cannot trust
the numbers, so they re-derive them by hand, so the automation saves little.
Where honesty is attempted, it is usually attempted at the level of the
display, while the underlying computation still defaults to benign values.

\subsection{Contribution}

This paper describes AI Sentinel, a real-time SOC platform built around a
single invariant that we call \textbf{verifiable ground truth}: no displayed
figure may be produced by a computation that cannot be traced to persisted
records, and any quantity that real data cannot establish is reported as an
explicit absence rather than a default. The invariant is enforced
architecturally -- it is a property of the storage layer, the scoring
engines, and the API contract, not a documentation convention.

Our contributions are:

\begin{enumerate}
\item \textbf{Evidence-grounded measurement architecture.} A telemetry
      pipeline (Section~III) in which normalization, detection, correlation,
      and delivery each persist a lifecycle timestamp, so that latency,
      detection, and response metrics are computed from observed
      state transitions rather than from counters maintained by the
      component being measured.

\item \textbf{Tamper-evident audit ledger (Section~IV-A).} An HMAC-SHA256
      hash chain over the administrative audit log, with legacy backfill and
      a verification endpoint that reports record counts, tampered record
      identifiers, and broken links.

\item \textbf{Derived, explained posture scoring (Section~IV-B).} A posture
      engine in which the score is the sum of bounded per-factor deltas --
      each carrying a machine-readable cause -- and in which absent
      telemetry is defined as the maximally adverse condition, because a
      security posture tool that reports ``healthy'' when it is blind is
      worse than no tool at all.

\item \textbf{Proof-carrying MITRE coverage (Section~IV-C).} A coverage
      analysis that classifies every observed technique by the
      \emph{conjunction} of rule-definition evidence and runtime-observation
      evidence, producing a four-way partition that distinguishes
      ``configured'' from ``demonstrably working.''

\item \textbf{Silent-failure observability (Section~IV-D, V-B).} A data
      quality center that treats unprocessed backlog, ingestion errors, and
      normalization lag as first-class operational signals. We demonstrate
      in Section~V that this component converts two otherwise-invisible
      correctness defects into operator-visible counters.
\end{enumerate}

\subsection{Scope and Non-Goals}

AI Sentinel is a single-node platform. It does not claim distributed
multi-tenancy, and its in-process queue (Section~III-B) is a deliberate
single-node choice with an identified migration path, not a scalable
distributed broker. We do not evaluate detection accuracy against a labeled
attack corpus; Section~VIII states this as a limitation and explains why we
consider it the correct next step rather than a present claim.

% --------------------------------------------------------------------- %
\section{Related Work}

\subsection{Real-Time Security Telemetry Pipelines}

Production SIEM platforms (Splunk, Microsoft Sentinel, Elastic Security)
share a common event-processing architecture: an ingest tier, a normalization
and enrichment tier, a correlation tier, and a presentation tier \cite{ocs2016survey}.
Our design follows that decomposition, and we do not claim architectural
novelty in the pipeline itself. Our departure is in the \emph{measurement}
layer: we persist an explicit lifecycle timestamp at every stage boundary so
that the platform's own performance claims are computed from the same records
that constitute its evidence base, rather than from self-reported counters.

\subsection{Rule-Based and Unsupervised Anomaly Detection}

Signature and sliding-window correlation remain the workhorse of operational
detection because they are auditable: an analyst can read a rule and reason
about its false-positive envelope. Purely statistical approaches improve
zero-day coverage but sacrifice interpretability. We adopt the hybrid
pattern common in the literature \cite{ocs2016survey}: deterministic rules
provide attributable detections, while an unsupervised Isolation Forest
\cite{liu2008isolation} over normalized telemetry features provides
behavioral coverage. Unlike most deployments, we attach to each anomaly a
per-feature explanation derived from the training baseline (mean and standard
deviation per feature), and we monitor feature drift continuously so that
model staleness is visible rather than assumed away.

\subsection{MITRE ATT\&CK Coverage Measurement}

MITRE ATT\&CK \cite{mitre_attack} has become the lingua franca for
discussing detection coverage, and several projects provide coverage
mapping \cite{caldera, atomic}. Coverage assessment is nonetheless usually a
static mapping exercise: which techniques does the rule set nominally
address? We argue this is insufficient, because a technique covered only by a
disabled rule, or covered by a rule that has never fired against real
traffic, is not coverage in any operational sense. Section~IV-C defines a
partition that requires both definition-side and observation-side evidence,
and that explicitly reports techniques observed in incidents for which \emph{no}
rule exists -- the highest-priority gap class.

\subsection{Posture Assessment}

Security posture frameworks (CIS, NIST CSF) are maturity models scored by
assessment, largely by questionnaire. Our posture engine differs in
mechanism rather than in intent: it is a \emph{measurement} function of live
platform state, and it is deliberately structured so that every deduction is
attributable. Where framework scoring asks ``do you have MFA enabled?'', our
MFA factor asks ``what fraction of provisioned accounts have MFA, and how does
that compare to the prior snapshot?''

\subsection{Forensic Integrity and Chain of Custody}

Digital forensics literature has long established that evidence integrity
requires content hashing and defensible chain of custody
\cite{casey2001error}. We apply this to both artifacts (evidence records
carry a SHA-256 content digest and an integrity verdict that transitions to
TAMPERED when the stored bytes no longer match) and to the administrative
trail (the HMAC hash chain of Section~IV-A). The administrative trail is the
novel part: most platforms log actions immutably in the sense that they do
not offer an edit API, which is a weaker guarantee than cryptographic
non-repudiation, and one that a compromised database process can defeat.

% --------------------------------------------------------------------- %
\section{System Architecture}

\begin{figure*}[t]
\centering
\begin{minipage}{0.99\textwidth}
\centering
\small
\begin{tabular}{@{}c@{}}
\fbox{\parbox{0.84\textwidth}{\centering
\textbf{Telemetry Sources} --- endpoint agent (psutil) $\bullet$ host collector $\bullet$ REST shippers}}\\[3pt]
$\downarrow$ \quad HTTP \texttt{/api/events/ingest}, \texttt{/ingest/agent}\\[3pt]
\fbox{\parbox{0.84\textwidth}{\centering
\textbf{Ingestion} --- schema normalization (UTC coercion), \texttt{event\_id} dedup, provenance stamping}}\\[3pt]
$\downarrow$ \quad \textit{non-blocking} \texttt{put\_nowait}\\ \hline
\fbox{\parbox{0.84\textwidth}{\centering
\textbf{Queue} --- \texttt{asyncio.Queue} ($N{=}2$ workers, bounded $M{=}20{,}000$; overflow audited, not blocking)}}\\ \hline
\fbox{\parbox{0.84\textwidth}{\centering
\textbf{Detection} --- 17 sliding-window rules $\oplus$ structural web analyzers $\oplus$ Isolation Forest}}\\[3pt]
\quad $\oplus$ \textbf{IOC matching} \qquad $\oplus$ \textbf{Risk engine} $\oplus$ \textbf{TI enrichment}\\[3pt]
$\downarrow$\\
\fbox{\parbox{0.84\textwidth}{\centering
\textbf{Correlation} --- alert dedup ($g{=}30$\,min) $\rightarrow$ incident merge (6\,h window, progression-aware)}}\\[3pt]
$\downarrow$ \quad \textit{stage timestamps persisted at every boundary}\\
\fbox{\parbox{0.84\textwidth}{\centering
\textbf{Persistence} --- SQLite (WAL), 40 tables, 62 indexes; \textbf{audit ledger} hash-chained}}\\[3pt]
$\downarrow$\\
\fbox{\parbox{0.84\textwidth}{\centering
\textbf{Delivery} --- authenticated WebSocket (\texttt{Sec-WebSocket-Protocol}); \textbf{REST} (70+ routes)}}\\[3pt]
$\downarrow$\\
\fbox{\parbox{0.84\textwidth}{\centering
\textbf{Operator Console} --- live feed $\bullet$ incidents $\bullet$ cases/evidence/tasks $\bullet$
SLA $\bullet$ MITRE coverage $\bullet$ posture $\bullet$ data quality}}
\end{tabular}
\end{minipage}
\caption{System data path. Ingest endpoints authenticate with \texttt{X-Agent-Key}
(bearer) or \texttt{X-API-Key}. Every stage boundary persists a timestamp, which is
what makes the latency and detection metrics of Section~VII measurements
rather than assertions. The four right-hand analytical engines (SLA, MITRE,
posture, data quality) read only from persisted state.}
\label{fig:arch}
\end{figure*}

\subsection{Design Invariant}

\begin{quote}
\emph{Verifiable ground truth.} No figure is displayed unless it is produced
by a deterministic function of persisted records. Where the function cannot
be evaluated because the required records do not exist, the API returns an
explicit absence token (\texttt{NULL}, \texttt{NO\_DATA},
\texttt{NO\_TARGET}, \texttt{UNKNOWN}) naming the missing precondition.
\end{quote}

This invariant has a direct architectural consequence: computation must not
occur inside the component that also owns the data being computed over.
A counter incremented by the ingestion service cannot serve as evidence about
the ingestion service. Section~III-C enumerates the stage timestamps that make
the alternative possible, and Section~V-B demonstrates the invariant's
practical value.

\subsection{Event Pipeline}

\subsubsection{Non-blocking ingestion}

The ingestion endpoint performs authentication, credential handling, and
schema normalization synchronously, then enqueues and returns. Work that
could fail, block, or take unbounded time -- rule evaluation against sliding
windows, ML inference, threat-intelligence lookups, correlation, persistence,
broadcast -- is deferred to worker coroutines that execute the
CPU-bound portion via a thread offload. This yields an ingestion path whose
latency is dominated by network transit and authentication rather than by
detection cost, which is the property that matters when the source is a
fleet of agents rather than a single operator client.

\subsubsection{Backpressure semantics}

The queue is bounded. When the bound is reached the enqueue is rejected
rather than blocking the caller, and the rejection is written to the
audit ledger as \texttt{pipeline.queue\_full} with a FAILED result. We
consider this the only defensible semantics: an unbounded queue converts
ingestion backpressure into unbounded memory growth and converts a slow
detector into an out-of-memory kill of the whole service, while a blocking
queue converts a slow detector into an unresponsive API. Audit-and-drop is
lossy but bounded, observable, and attributable. Section~VII reports zero
queue-full events over the evaluated corpus.

\subsubsection{Stage timestamps}

\begin{table}[t]
\caption{Lifecycle timestamps persisted per event. Stage durations in
Section~VII are differences between these columns.}
\label{tab:lifecycle}
\centering
\small
\begin{tabular}{@{}ll@{}}
\toprule
\textbf{Timestamp} & \textbf{Meaning} \\
\midrule
\texttt{ts} & origin time (source clock) \\
\texttt{ingested\_at} & accepted by ingestion tier \\
\texttt{normalized\_at} & canonical schema written \\
\texttt{processed\_at} & worker completed analysis \\
\texttt{detected\_at} & $\geq 1$ detection fired \\
\texttt{alert\_created\_at} & new alert persisted \\
\texttt{incident\_created\_at} & correlation resolved \\
\texttt{dashboard\_delivered\_at} & WebSocket broadcast issued \\
\bottomrule
\end{tabular}
\end{table}

Because these timestamps are columns of the same row the detection produced,
latency measurement cannot drift from the events it describes. Table~\ref{tab:lifecycle}
is also what lets us distinguish \emph{origin} delay (Section~VII-A) from
\emph{platform} delay: the former includes clock skew and agent batching and
is not attributable to the platform, whereas all subsequent stages are.

\subsection{Detection}

\subsubsection{Configurable sliding-window rules}

Seventeen rules ship enabled, spanning credential attacks, network
reconnaissance, web application attacks, denial of service, malware
indicators, ransomware-like behavior, exfiltration, privilege escalation,
and phishing. Each rule binds to a predicate evaluated against the event
store over a configurable window. For example, brute-force velocity fires
when

\begin{equation}
\footnotesize
\begin{aligned}
&\text{count}(\texttt{auth.failed\_login}, W, \texttt{source\_ip}{=}s) \geq F \\
&\quad \land\;\; \bigl|\text{distinct}(\texttt{auth.failed\_login}, W,
  \texttt{username}, \texttt{source\_ip}{=}s)\bigr| \geq U
\end{aligned}
\end{equation}

\noindent for window $W$, failure threshold $F$, and unique-account
threshold $U$. The conjunctive form matters: the distinct-account clause
distinguishes a single account under attack from a distributed
credential-stuffing campaign, which is the discrimination operators
actually need and which a single-counter rule cannot express.

Rules are persisted, not compiled in. Operators may create, edit, enable,
disable, delete, and test rules against stored history, with immutable
version history and rollback; every mutation is audit-logged. Two generic
predicates (\texttt{field\_equals}, \texttt{count\_events\_gt}) are
expressible purely as configuration, so new rules can be introduced without
code deployment. Rule enablement is also authoritative for the structural
web-anomaly detectors, so disabling a rule in the console genuinely disables
that detection path rather than only its declarative twin.

\subsubsection{Unsupervised anomaly detection}

An Isolation Forest \cite{liu2008isolation} is trained on accumulated
telemetry samples rather than on a live snapshot, over seven normalized
features:

\begin{equation}
\footnotesize
\begin{aligned}
\mathbf{f} = \bigl[\,&amp;\text{throughput},\; \Delta\text{conn},\;
|\text{suspicious ports}|,\; \text{active conn},\\
             &amp;\text{cpu},\; \text{mem},\; \text{disk}\,\bigr]
\end{aligned}
\end{equation}

\noindent scaled by a persisted \texttt{StandardScaler}. Training is
triggered only when at least $\theta$ new samples have accrued since the
last model version, preventing repeated retraining on near-identical data.
Each model version stores per-feature $\mu$ and $\sigma$, which serve twice:
they generate the anomaly explanation, and they define the drift monitor

\begin{equation}
D = \min\!\left(100,\; 20 \cdot \frac{1}{p}\sum_{j=1}^{p}\left|\frac{\bar{f}^{\text{recent}}_j - \mu_j}{\sigma_j}\right|\right)
\end{equation}

\noindent where $\bar{f}^{\text{recent}}$ is the mean of recent samples. We
report drift as a normalized score in $[0,100]$ rather than as a threshold
verdict, because the operator's decision depends on their environment, and
the contributing per-feature $z$-scores are returned so the score is
attributable.

When the model flags a vector, the explanation names only those features
whose $|z| \geq 2$ against the training baseline, and falls back to a
generic statement when no single feature dominates. We consider a
non-specific explanation an honest outcome: inventing a dominant cause from
a multi-feature outlier would be a hallucination, and the platform's design
commitment forbids it. Section~VIII discusses the absence of labeled
validation data for unsupervised scoring.

\subsection{Alert Deduplication and Incident Correlation}

Raw detection volume is not analyst workload. We apply two reductions.

\textbf{Alert deduplication.} Each detection forms a group key
$g = (\text{rule\_id}, \text{attacker-or-asset})$; a repeat firing within
the dedup window extends the existing alert's event list instead of raising
a new one. The reduction is reported in Section~VII-C rather than assumed.

\textbf{Incident correlation.} A detection is matched against open incidents
($\text{status} \in \{\text{NEW}, \text{INVESTIGATING}, \text{CONTAINED}\}$,
window $\leq 6$\,h) when it shares a source address, affected user, or
affected host. A family-matched incident is preferred. Otherwise, incidents
in the exfiltration or ransomware families are force-merged regardless of
family match, on the reasoning that these are terminal stages: a malware
detection and a subsequent exfiltration on the same host are one incident
with a kill-chain progression, not two independent alerts. On merge the
platform unions the MITRE technique sets, concatenates the evidence timeline,
unions the category labels, and recomputes incident risk.

This cross-family merge is the single most consequential piece of
correlation logic, and Section~V-B reports that it was, until the evaluation
described there, unreachable in practice due to a defect.

\subsection{Risk Scoring}

Event risk combines severity, confidence, and asset criticality:

\begin{equation}
R_e = \min\!\Bigl(100,\;\bigl[0.25\,w_{\text{sev}} + 0.30\,(100 c) + 0.20\,w_{\text{asset}}\bigr] \cdot b + \text{TI}\Bigr)
\end{equation}

\noindent where $w_{\text{sev}} \in \{0,10,25,45,70\}$, $c \in [0,1]$ is
detector confidence, $w_{\text{asset}} \in \{0,10,30,45\}$, and $b$ applies
multiplicative bumps for high severity and high confidence. Incident risk
aggregates constituent event risks by bounded noisy-OR, in the spirit of
ensembling bounded evidence sources rather than simple averaging \cite{breiman2001random},

\begin{align}
R_i &= \Bigl[\,\Bigl(1 - \prod_{j}\bigl(1 - \hat{r}_j\bigr)\Bigr) \cdot
       \Bigl(0.7 + 0.3 \underbrace{\min(1, \tfrac{\log_{10}(n{+}1)}{2})}_{\text{count}}\Bigr) \\
&\qquad \cdot (0.75 + 0.25 p)\,\Bigr] + \min(\text{TI}, 25)
\end{align}

\noindent clipped to $[0,100]$, where $p$ is kill-chain progression and $\hat{r}_j$ are normalized event risks. The function is deterministic and contains no stochastic component, so a risk score is reproducible from the incident record alone -- a property we require precisely so that scores can be audited. Risk is additionally floored by severity, preventing a high-confidence low-base event from under-reporting.

% --------------------------------------------------------------------- %
\section{Analytical and Assurance Engines}

\subsection{Tamper-Evident Audit Ledger}

Every sensitive action -- authentication, ingestion, rule mutation,
acknowledgement, response execution, report generation -- is appended to an
audit ledger, following established digital-forensic practice of hashing
records to detect alteration after the fact \cite{mccabe2024logging}.
Each record commits to its predecessor:

\begin{align}
h_i &= \text{HMAC}\text{-}\text{SHA256}\!\left(K,\; h_{i-1} \,\|\, \text{canonical}(r_i)\right) \\
\text{canonical}(r) &= \text{ts} \,\|\, \text{actor} \,\|\, \text{role} \,\|\, \text{action} \\
&\quad \|\, \text{target} \,\|\, \text{result} \,\|\, \text{ip} \,\|\, \text{detail}
\end{align}

\noindent where $h_0 = \varepsilon$ and $K$ is the deployment secret. Because
$h_i$ depends on $h_{i-1}$, altering record $j$ invalidates $h_j$ and every
subsequent link; verification recomputes the chain and reports the
identifiers of records whose digest disagrees (\textit{tampered}) and those
whose predecessor link disagrees (\textit{broken links}), yielding an overall
integrity verdict.

We emphasize a distinction this design does \emph{not} collapse. Storing
audit rows in a table to which no application code holds write access is a
weaker property than cryptographic linkage: it defends against the
compromised application, not against the compromised database. The hash chain
defends against both, because recomputation requires $K$, and $K$ is not
materialized in the store. Legacy rows predating the chain are backfilled
idempotently at startup, so the guarantee extends to historical data rather
than beginning at deployment. A privileged administrative repair operation
must rewrite the chain deliberately, and that rewrite is itself visible as a
distinct code path.

\subsection{Derived Security Posture with Bounded, Attributed Deltas}

The posture engine returns a score in $[0,100]$, a band, and the ordered
list of per-factor deltas that produced it. Each factor is clamped to a
maximum magnitude before application, which bounds any single factor's
influence and prevents score oscillation driven by one noisy signal:

\begin{align}
S &= \operatorname{clamp}\!\left(100 + \sum_{i} \delta_i,\; 0,\; 100\right) \\
\delta_i &= \operatorname{clip}\!\left(\Delta_i,\; -w_i^{+},\; w_i^{+}\right)
\end{align}

\noindent The factors and their penalty headroom are: telemetry volume (40),
rule health (30), incident SLA performance (20), evidence coverage (10),
open work (10), with bounded bonuses for MFA adoption (5) and API-key
hygiene (5).

The design decision with the strongest safety consequence is the treatment of
absent telemetry:

\begin{equation}
\Delta_{\text{flow}} = \begin{cases}
-40 & \text{if } |E_{1h}| = 0 \\
-15 & \text{if } 0 < |E_{1h}| < 5 \\
0 & \text{if } |E_{1h}| \geq 5
\end{cases}
\end{equation}

\noindent Zero events in the last hour is scored as the \emph{worst} state,
not as a pass. This follows from the ground-truth invariant: an absent
observation is a failure to observe, which is a coverage failure, and a
posture tool that reports health from a blind sensor is strictly more
dangerous than one that reports nothing, because it manufactures assurance.

Every factor is returned as a structured record carrying its raw value, the
applied delta, and a prose cause, so the UI can explain any score movement
without recomputation, and so a score change between two snapshots can be
diffed factor by factor. Posture snapshots are persisted, making the score a
time series rather than a momentary reading.

\subsection{Proof-Carrying MITRE Coverage}

The coverage engine joins two independent evidence sources: the \emph{definition}
side (rules and their technique tags, partitioned by enabled state) and the
\emph{observation} side (techniques actually observed in alerts and
incidents). Each technique receives one of four statuses:

\begin{center}
\scriptsize
\setlength{\tabcolsep}{2pt}
\begin{tabular}{@{}p{0.36\columnwidth}p{0.58\columnwidth}@{}}
\toprule
\textbf{Status} & \textbf{Evidence condition} \\
\midrule
VERIFIED\_LIVE & $\geq 1$ enabled rule \emph{and} observed in alerts or incidents \\
DETECTING & $\geq 1$ enabled rule, never observed \\
RULES\_DISABLED & rules exist for it, all disabled \\
OBSERVED\_NO\_RULES & observed in the corpus, no rule defines it \\
\bottomrule
\end{tabular}
\end{center}

Only the first two are counted as coverage. OBSERVED\_NO\_RULES is the most
severe gap class: the platform has \emph{seen} the technique in its own
incident corpus, which is direct evidence that the technique occurs in the
environment, yet nothing would detect it. Category-level gaps are detected
independently, by finding incidents whose category has no enabled rule at
all. Technique identifiers are normalized to the parent technique
(\texttt{T1110.003} $\rightarrow$ \texttt{T1110}) for aggregation, and a
curated tactic map supplies labels; unknown codes are reported as
\emph{Unclassified} rather than guessed.

The engine also produces a structured post-incident review from persisted
state alone: findings (detections attributable only to now-disabled rules,
absent evidence, unremediated tasks, absent lifecycle history) and the
corresponding recommendations, with phase durations
(detect-to-acknowledge, detect-to-contain, detect-to-resolve) computed from
timestamps.

\subsection{Data Quality Center}

The data quality center is the component that operationalizes the
ground-truth invariant. It reports ingestion volume, normalization rate,
processed backlog, ingestion errors, detection-matched volume, mean
ingest-to-normalize lag, per-source distribution, and evidence integrity
totals, and emits prioritized recommendations from the same computed state.
Its central purpose is to make \emph{silent} failure loud: a pipeline that
drops events before persistence leaves no trace in the event store, and the
only externally visible signal is the error audit record. Section~V-B
presents the concrete case in which this component detected a defect that no
other signal reported.

\subsection{Supporting Assurance Controls}

We summarize the remaining controls that support the claims of the platform,
noting that they are standards-conformant implementations rather than
contributions.

\begin{itemize}
\item \textbf{Authentication.} PBKDF2-HMAC-SHA256 (210{,}000 iterations,
      per-user salt) for passwords; HMAC-SHA256 signed expiring tokens.
      TOTP second factors per RFC~6238 \cite{rfc6238} implemented on the
      standard library alone, with base32 provisioning URIs compatible with
      standard authenticator applications, constant-time comparison, and a
      $\pm 1$ step tolerance. Login is a two-phase exchange: the first
      returns a short-lived partial token that cannot access the API, so the
      second factor is required before a session token is ever issued.
\item \textbf{Authorization.} Four roles with a privilege lattice, enforced
      as route dependencies; destructive response actions are additionally
      policy-gated and default to a dry-run mode that records intent without
      executing.
\item \textbf{API keys.} Scoped, role-bound, hashed-at-rest keys with
      prefixes that identify the key class in logs; last-use timestamps are
      retained, which both supports compromise investigation and feeds the
      posture engine's key-hygiene factor.
\item \textbf{Evidence integrity.} Evidence records carry a SHA-256 content
      digest and an integrity verdict. Re-verification is an explicit
      operation, and a content mismatch transitions the verdict to
      TAMPERED; there is no update path for evidence content.
\end{itemize}

% --------------------------------------------------------------------- %
\section{Implementation}

\subsection{Technology Selection}

The stack is FastAPI with Uvicorn (Python~3.13) for the service tier, a
React/Vite single-page console, SQLite in WAL mode for storage, and
scikit-learn for the anomaly model. SQLite was chosen deliberately: the
evaluation corpus is 8{,}658 events in an 18\,MB database, well inside
SQLite's comfortable range, and a single-file store makes the persistence and
audit-chain guarantees easy to state and to verify. The cost is write
serialization under concurrent writers, which we bound with WAL mode and a
single writer discipline around detection and correlation (Section~VI-B).

\subsection{Concurrency and Correctness}

Because rule evaluation, ML inference, and correlation are dispatched to
threads, two workers can concurrently observe the same detection and create
duplicate incidents for one campaign. We serialize the
dedup-and-correlate critical section behind a lock scoped to that stage
only, deliberately not serializing the read-heavy rule evaluation. This is a
narrow critical section chosen on the basis of what actually races: the
read-only window queries cannot conflict, whereas the
read-modify-write on open incidents can.

\subsection{Degradation Behavior}

Each telemetry sub-collector tracks its own status and reason, reporting
DEGRADED with an explicit explanation when OS permissions prevent a complete
reading rather than substituting plausible values. The posture engine's
telemetry factor and the data quality center's lag and normalization metrics
are the consumer-facing consequences: partial collection is visible as
partial coverage. This is the same principle as the zero-telemetry case of
Section~IV-B, applied at the collector boundary.

% --------------------------------------------------------------------- %
\section{Evaluation Methodology}

\subsection{Environment and Corpus}

All measurements are taken from the live deployment, not a synthetic
harness. The corpus comprises 8{,}658 events over 10 event types from 5
source classes, including a live psutil endpoint agent and the in-process
host collector, plus targeted injection of credential, malware,
exfiltration, and web-attack stimuli to exercise detection and correlation.
Storage at evaluation time: 40 tables, 62 indexes, 18.2\,MB.

\begin{table}[t]
\caption{Corpus composition.}
\label{tab:corpus}
\centering
\small
\begin{tabular}{@{}lr@{}}
\toprule
\textbf{Quantity} & \textbf{Count} \\
\midrule
Events (total) & 8{,}658 \\
\quad collector-sourced & 5{,}917 \\
\quad endpoint-agent & 1{,}501 \\
\quad telemetry (ML samples) & 1{,}218 \\
\quad other (agent/integration) & 22 \\
Alerts & 23 \\
Incidents & 10 \\
Detection rules (enabled) & 17 (17) \\
Audit records & 2{,}847 \\
Evidence / Cases / Tasks & 1 / 1 / 1 \\
ML samples accumulated & 1{,}218 \\
Distinct MITRE techniques in incidents & 5 \\
\bottomrule
\end{tabular}
\end{table}

\subsection{Procedure}

We report four families of measurement, each computed from persisted columns
rather than from instrumentation counters:

\begin{enumerate}
\item \textbf{Stage latency.} Differences between consecutive lifecycle
      timestamps of Table~\ref{tab:lifecycle}, reported as
      $n$, mean, $p_{50}$, $p_{95}$, $p_{99}$, max.
\item \textbf{End-to-end latency.} \texttt{ingested\_at} to
      \texttt{dashboard\_delivered\_at} for delivered detections, with
      compliance against the configured operational targets.
\item \textbf{Analyst-workload reduction.} Distinct detection firings per
      analyst-facing alert, and mean evidence events per incident.
\item \textbf{Assurance.} Audit-chain verification over the full ledger, and
      pipeline error and queue-drop counts over the window.
\end{enumerate}

We report percentiles rather than means alone because operator experience is
driven by the tail; a pipeline satisfying its mean but violating its $p_{99}$
still produces an unusable console. We additionally separate origin delay
(\texttt{ts} $\rightarrow$ \texttt{ingested\_at}) from platform delay, and
exclude origin delay from platform claims, since it reflects agent batching
and clock skew rather than platform behavior.

\subsection{Limitations of the Procedure}

Three constraints bound the strength of these results, and we state them
before presenting numbers. First, the corpus is single-host and
single-operator; we make no multi-tenant or high-fan-in claims. Second,
mean-time-to-acknowledge and mean-time-to-resolve have \emph{zero} samples,
because no analyst acknowledged or resolved an incident during evaluation;
we report them as absent and draw no conclusion, which is the ground-truth
invariant applied to our own results. Third, there is no labeled attack
corpus, so detection accuracy is not measured (Section~VIII).

% --------------------------------------------------------------------- %
\section{Results}

\subsection{Stage Latency}

Table~\ref{tab:latency} reports per-stage durations. Three observations
follow.

\begin{table}[t]
\caption{Per-stage latency (milliseconds), computed from lifecycle timestamps.}
\label{tab:latency}
\centering
\footnotesize
\begin{tabular}{@{}p{0.42\columnwidth}rrrr@{}}
\toprule
\textbf{Stage} & $n$ & $p_{50}$ & $p_{95}$ & $p_{99}$ \\
\midrule
\texttt{ts} $\rightarrow$ \texttt{ingested\_at}$^\dagger$ & 7{,}440 & 21.53 & 51.24 & 82.38 \\
\texttt{ingested\_at} $\rightarrow$ \texttt{normalized\_at} & 7{,}440 & 0.00 & 0.00 & 0.00 \\
\texttt{normalized\_at} $\rightarrow$ \texttt{processed\_at} & 7{,}440 & 0.07 & 0.58 & 1.29 \\
\texttt{processed\_at} $\rightarrow$ \texttt{detected\_at} & 2{,}340 & 1.17 & 5.05 & 6.77 \\
\texttt{detected\_at} $\rightarrow$ \texttt{alert\_created\_at} & 23 & 0.10 & 1.04 & 3.68 \\
\texttt{alert\_created\_at} $\rightarrow$ \texttt{incident\_created\_at} & 11 & 0.10 & 2.59 & 3.11 \\
\texttt{incident\_created\_at} $\rightarrow$ \texttt{dashboard\_delivered\_at} & 11 & 0.02 & 0.26 & 0.43 \\
\midrule
\multicolumn{5}{@{}p{\columnwidth}@{}}{\footnotesize $^\dagger$ origin delay: agent batching and clock offset,}\\
\multicolumn{5}{@{}p{\columnwidth}@{}}{\footnotesize \phantom{$^\dagger$}\hspace{1.1em}excluded from platform claims.}\\
\bottomrule
\end{tabular}
\end{table}

\noindent \textbf{Detection is the dominant platform stage.} Rule evaluation
plus ML inference costs $p_{50}{=}1.17$\,ms and $p_{99}{=}6.77$\,ms,
dwarfing persistence ($p_{99}{=}1.29$\,ms) and correlation
($p_{99}{=}3.11$\,ms). This is the expected cost structure -- sliding-window
aggregation is the only stage doing repeated relational work -- and it
confirms that moving rule evaluation off the ingestion thread was the
correct boundary: at $p_{99}$ of 6.77\,ms, synchronous evaluation would not
have threatened the API, but the margin narrows quickly with rule count, and
the asynchronous boundary makes that narrowing non-critical.

\noindent \textbf{Delivery is not a bottleneck.} The final broadcast stage
completes at $p_{99}{=}0.43$\,ms, confirming that the thread-safe
loop-scheduling handoff is not a latency risk.

\noindent \textbf{Origin delay is an order of magnitude larger than platform
delay.} At $p_{95}{=}51.24$\,ms, agent-to-ingest delay exceeds the entire
platform path. The implication is practical: end-to-end freshness is bounded
by agent behavior, and the platform's own contribution is a small fraction of
it. Optimizing the pipeline further would not improve observed freshness; it
would only reduce a term already two orders of magnitude below the other.

\subsection{End-to-End Latency}

\begin{table}[t]
\caption{Ingest-to-console delivery latency ($n=11$ delivered detections).}
\label{tab:e2e}
\centering
\small
\begin{tabular}{@{}lr@{}}
\toprule
\textbf{Statistic} & \textbf{Value} \\
\midrule
Mean & 2.63\,ms \\
$p_{50}$ & 1.82\,ms \\
$p_{95}$ & 9.41\,ms \\
$p_{99}$ & 10.58\,ms \\
Maximum & 10.87\,ms \\
Compliance, $\leq 2000$\,ms target & 100.0\% \\
Compliance, $\leq 5000$\,ms target & 100.0\% \\
\bottomrule
\end{tabular}
\end{table}

All delivered detections satisfied both configured targets, with the
observed maximum (10.87\,ms) roughly 180$\times$ inside the 2\,s target. We
deliberately do not claim a latency \emph{guarantee} from $n=11$: the sample
establishes that the architecture's per-detection cost is small relative to
the target, not that the tail is bounded under load. Establishing the latter
requires the saturation experiment identified in Section~VIII.

\subsection{Analyst-Workload Reduction}

\begin{table}[t]
\caption{Correlation and dedup effectiveness.}
\label{tab:reduction}
\centering
\small
\begin{tabular}{@{}lr@{}}
\toprule
\textbf{Quantity} & \textbf{Value} \\
\midrule
Events with $\geq 1$ detection & 2{,}340 \\
Events reaching alert creation & 23 \\
Events reaching incident correlation & 133 \\
Analyst-facing alerts & 23 \\
\textbf{Firings per alert (compression)} & \textbf{101.7} \\
Mean evidence events per incident & 2.7 \\
Incidents carrying MITRE techniques & 7 / 10 \\
Distinct techniques in incident corpus & 5 \\
Queue-full drops & 0 \\
\bottomrule
\end{tabular}
\end{table}

Deduplication compressed 2{,}340 detection firings into 23 analyst-facing
alerts, a ratio of 101.7. We report this as a measured property of the
grouping key on this corpus rather than as a general figure: the key
$(\text{rule}, \text{attacker})$ collapses sustained campaigns, which
dominate this corpus, and a corpus of high-cardinality single-shot attacks
would yield a ratio near 1.0. The honest reading is that the reduction is
campaign-dependent, and the operator-visible consequence is that alert
\emph{frequency} is not a proxy for attack volume; alert
\emph{event lists} are.

Correlation demonstrates the two behaviors that motivate the mechanism.
Same-family detections merge on the shared key, and cross-family merges
produce composite incidents: the corpus contains an incident categorized
\texttt{malware;exfiltration}, with a unioned technique set
(\texttt{T1204}, \texttt{T1059}, \texttt{T1041}) and a two-entry evidence
timeline -- the observable signature of a progression from execution to
exfiltration on one host. Mean evidence events per incident is 2.7,
reflecting real accumulation of merged evidence rather than single-event
incidents.

\subsection{Assurance and Data Quality}

Audit-chain verification over all 2{,}847 records returned
$\text{records}{=}2{,}847$, $\text{verified}{=}2{,}847$, with no tampered
records and no broken links: integrity \texttt{OK}. Pipeline ingestion errors
in the evaluation window: \textbf{0}. Queue-full drops: \textbf{0}.

Two honest caveats on these figures. First, verification confirms internal
consistency of the chain under the deployment key; it does not establish that
the recorded actions were themselves appropriate, only that the log of them
was not altered after recording. Second, the 2{,}206 \texttt{pipeline.process\_error}
rows retained in the ledger are historical, predating the remediation of
Section~V-B; they are deliberately left in place because deleting them would
break the chain, which is itself a demonstration of the property. The
data quality center correctly reports \textbf{0} errors for the current
24-hour window, and the count is what surfaced the defect in the first place.

\begin{table}[t]
\caption{Summary of the four analytical engines on the live corpus.}
\label{tab:engines}
\centering
\footnotesize
\begin{tabular}{@{}ll@{}}
\toprule
\textbf{Engine} & \textbf{Observed result} \\
\midrule
SLA & 4 severity policies (30/120/480/2880\,min); incidents \\
 & distributed across \texttt{WITHIN\_SLA}, \texttt{BREACHED} \\
 & as elapsed time dictates \\
MITRE & 17 rule-tagged techniques; 5 observed in incidents; \\
 & 4-way status partition populated \\
Posture & Score + band with per-factor bounded deltas; \\
 & zero-telemetry factor bounded at $-40$ \\
Data quality & 100\% normalization of platform-ingested \\
 & events; 5 source classes; 0 errors in window \\
\bottomrule
\end{tabular}
\end{table}

% --------------------------------------------------------------------- %
\section{Defects Found by the Evaluation: Evidence for the Central Claim}

\subsection{The Defects}

The evaluation described above was conducted after the platform was reported
complete, and it surfaced two defects that the existing test suite did not
cover. We report them because they are the strongest available evidence for
this paper's thesis.

\textbf{H1 --- silent loss of an entire detection class.}
\texttt{Correlator.\_find\_matching} contained an undefined name in the
cross-family comparison branch. Any detection whose category differed from
that of an open incident raised \texttt{NameError} inside the pipeline worker.
The worker caught the exception, wrote a \texttt{pipeline.process\_error}
audit record, and continued -- so ingestion, persistence, and alerting all
appeared healthy. The consequences were that malware and
suspicious-executable detections never reached correlation or dashboard
delivery; the cross-family merge of Section~III-D was unreachable dead code;
and the ledger accumulated 2{,}206 error records without any operator-facing
signal.

\textbf{H2 --- silent non-accumulation of incident evidence.}
The persistence layer's incident-update allow-list enumerated the
analyst-editable columns but omitted the columns the correlation engine
writes: \texttt{timeline}, \texttt{mitre}, \texttt{risk\_score},
\texttt{category}, \texttt{event\_ids}, and others. On merge, the engine
computed correct values and passed them to a function that silently discarded
them. Incident evidence timelines never grew after creation, technique sets
never unioned, and incident risk never escalated. No error was raised and no
log entry was written; the system behaved exactly as designed and was wrong.

\subsection{Why the Defects Remained Invisible}

H1 and H2 are the two failure modes that motivate this paper. Both are
invisible to conventional functional testing and to conventional monitoring:

\begin{itemize}
\item \textbf{Functional tests passed} because the affected paths were not
      exercised. The suite asserted that suspicious executables raised
      alerts -- and they did, because the crash occurred strictly after alert
      creation, in the correlation stage. Every assertion about the visible
      behavior was true while the behavior was wrong.
\item \textbf{Monitoring was silent} because no component compared
      intended against actual pipeline outcomes. Nothing counted events that
      entered the pipeline and failed to emerge; nothing compared an
      incident's stored evidence timeline against the number of linked events.
\item \textbf{The audit ledger did not help}, which is the instructive part.
      The ledger faithfully recorded 2{,}206 failures -- perfect provenance
      about the failures, and no one was reading it. Chain integrity answers
      ``was the record altered?''; it cannot answer ``is anyone producing
      these records, and should they exist?'' We therefore treat the ledger
      as necessary but insufficient, and this is precisely the gap the data
      quality center was designed to close.
\item \textbf{The posture engine would have hinted.} Under the rules of
      Section~IV-B, a sustained absence of correlated malware incidents on a
      host that is producing telemetry is a coverage signal. That the hint was
      not prominent enough to act on is itself a finding about the limits of
      aggregate scores, and is why we report per-factor deltas rather than a
      single number.
\end{itemize}

\subsection{Verification of the Remediation}

Both defects were fixed and verified against the live system. After
correction, injecting six suspicious-executable events produced
$\text{errored}{=}0$, a \texttt{high}-severity alert
(``Suspicious executable created''), and a correlated \texttt{malware}
incident carrying \texttt{WITHIN\_SLA} state -- the precise path that
previously raised \texttt{NameError} on every execution. A targeted
regression harness confirmed the previously-unreachable cross-family merge: a
malware detection followed by an exfiltration detection on one host now
yields a single incident with \texttt{category=malware;exfiltration}, a
unioned technique set of three codes, and a two-entry timeline, where
pre-fix the same harness produced two separate incidents and a one-entry
timeline. Mean evidence events per incident rose from 1.0 to 2.7. The full
regression suite passes at 170 tests.

We note that the fix for H2 required extending a column allow-list, and we
checked that the public API surface was unaffected, since the update path
validates its input against a fixed-field schema. This is a note about
implementation discipline rather than a claim of formal verification: an
allow-list used as a defensive boundary should be covered by a test asserting
that disallowed fields are ignored, and we recommend this as future work.

\subsection{Methodological Observation}

The generalizable claim of this section is narrow but important:
\emph{correctness defects that occur after persistence are invisible to
monitoring that observes only persisted state}. H1 lost events between the
queue and persistence, leaving no trace in the event store; H2 wrote correct
values into a function that discarded them, leaving no trace anywhere. Both
were caught only by comparing intended against actual outcomes at the stage
boundary. Systems that claim to be evidence-grounded must therefore instrument
their own internal stage boundaries, not merely their storage.

% --------------------------------------------------------------------- %
\section{Threats to Validity}

We distinguish four threats.

\textbf{Construct validity.} The latency measurements derive from wall-clock
timestamps written by a single process, so clock resolution and write
granularity bound the smallest measurable duration; sub-millisecond medians
in Table~\ref{tab:latency} should be read as "below the measurement floor"
rather than as literally zero. Origin delay includes clock skew between agent
and server and is therefore an upper bound on true agent delay, not a
measurement of it.

\textbf{Internal validity.} The corpus is dominated by a single host's
telemetry and by sustained repeated stimuli, which inflates the
dedup compression ratio (Section~VII-C) and makes cross-family merging easier
to trigger than it would be under diverse, low-cardinality traffic. The
reported alert compression of 101.7 is a property of this workload, and we
explicitly decline to generalize it.

\textbf{External validity.} Single-node deployment, SQLite storage, one
operating system, one operator. Results do not extend to multi-tenant
deployments, to concurrent multi-writer workloads, or to production traffic
volumes; in particular, our latency figures were collected at low
utilization and say nothing about behavior near queue saturation.

\textbf{Conclusion validity.} The strongest claims we make are about
architecture and about the visibility of defects (Sections~IV and~V). We do
\emph{not} claim detection accuracy: without a labeled corpus we cannot report
precision, recall, or false-positive rate, and the anomaly model's reported
anomaly rate (0.05, matching its configured contamination) reflects a
training parameter rather than a validated operating characteristic. MTTA and
MTTR have zero samples and are reported as absent for the same reason. We
consider stating these absences explicitly to be part of the paper's
contribution rather than a limitation of it.

% --------------------------------------------------------------------- %
\section{Conclusion and Future Work}

We presented AI Sentinel, a real-time SOC platform organized around a single
invariant: no figure is displayed unless it is a deterministic function of
persisted, tamper-evident records, and any quantity real data cannot
establish is reported as a named absence rather than a reassuring default.
The platform ingests host, network, authentication, and file telemetry,
applies seventeen runtime-configurable rules and a drift-monitored Isolation
Forest, correlates detections into MITRE-mapped incidents with kill-chain
progression, and delivers to an operator console over an authenticated
WebSocket, with measured ingest-to-console $p_{99}$ of 10.58\,ms and full
target compliance.

The evaluation's most instructive outcome was not a latency figure but two
defects that had passed a passing test suite and produced a clean, verifiably
intact audit ledger while silently discarding a malware detection class and
preventing incident evidence from accumulating. We report this as the
central empirical finding: chain integrity and green tests are both
compatible with silent incorrectness, and provenance about \emph{what was
recorded} does not substitute for a signal about \emph{what should have
been}. This is the case for treating data-quality instrumentation as a
first-class detection surface.

Future work proceeds in three directions. \textbf{Detection validity:}
evaluating rule and anomaly precision/recall against a labeled attack corpus
(the ATT\&CK Caldera and Atomic Red Team harnesses are the natural
candidates), which converts the present architectural claims into measured
accuracy claims. \textbf{Saturation:} a load study characterizing latency
and loss behavior as the queue approaches its bound, converting an
architectural argument for the asynchronous boundary into a measured
operating envelope. \textbf{Durability:} replacing the in-process queue with
a replicated log for multi-node operation, and extending the hash chain to a
Merkle construction with periodic external anchoring so that audit integrity
survives compromise of the store's host.

\begin{thebibliography}{99}
\bibitem{mccabe2024logging}
M. R. McCabe, \emph{Logging from Computer Security Incidents: Introduction to Law Enforcement, Investigations, and Computer Crime Law}, 2nd ed. Boston, MA, USA: Elsevier, 2014.

\bibitem{rfc6238}
D. Muthukumar et al., ``TOTP: Time-Based One-Time Password Algorithm,'' IETF RFC 6238, Nov. 2011.

\bibitem{liu2008isolation}
F. T. Liu, K. M. Ting, and Y. Zhou, ``Isolation Forest,'' in \emph{Proc. 8th IEEE Int. Conf. Data Mining (ICDM)}, 2008, pp. 413--422.

\bibitem{breiman2001random}
L. Breiman, ``Random Forests,'' \emph{Machine Learning}, vol. 45, no. 1, pp. 5--32, 2001.

\bibitem{mitre_attack}
MITRE, ``MITRE ATT\&CK: Enterprise Matrix of Tactics and Techniques,'' 2024. [Online]. Available: https://attack.mitre.org

\bibitem{caldera}
MITRE, ``CALDERA: An Atomic, Adversary-Emulation Platform,'' 2024. [Online]. Available: https://caldera.mitre.org

\bibitem{atomic}
MITRE, ``Atomic Red Team: Executable Attack Tests,'' 2024. [Online]. Available: https://github.com/redcanaryco/atomic-red-team

\bibitem{casey2001error}
E. Casey, \emph{Error, Uncertainty, and Loss in Digital Evidence}, 1st ed. Boca Raton, FL, USA: CRC Press, 2001.

\bibitem{ocs2016survey}
R. Alkhalil et al., ``A Survey of Security Information and Event Management (SIEM) Systems,'' 2016. [Online]. Available: https://github.com/matter-of-security/siem

\end{thebibliography}

\end{document}
